\section*{अध्याय \ref{ch:b3:statistical-thermo-applied} --- सांख्यिकीय ऊष्मागतिकी के अनुप्रयोग}

\begin{solution}{exo:b3:statistical-thermo-applied:1}
$\frac32R + R + (3 \times 3 - 5)R = 6.5R = \qty{54.0}{J.K^{-1}.mol^{-1}}$। कक्ष ताप पर स्थानांतरण
और घूर्णन चिरसम्मत हैं, दोनों तनन जमे हुए हैं और द्वि-अपभ्रष्ट
बंकन, जो सबसे निम्न कंपन है, आंशिक रूप से उत्तेजित है।
\end{solution}

\begin{solution}{exo:b3:statistical-thermo-applied:2}
$x = 1$: $\eu/(\eu - 1)^2 = 0.921$। $x = 10$: $100\eu^{10}/(\eu^{10} - 1)^2 \approx 100\eu^{-10} = \num{4.5e-3}$।
\end{solution}

\begin{solution}{exo:b3:statistical-thermo-applied:3}
\qty{20}{K}: $x_{\mathrm{para}} = 1/(1 + 3 \times 3\eu^{-2 \times 85.35/20}) = 1/(1 + 9 \times \num{1.96e-4}) = 0.998$।
\qty{77}{K}: सम योग $1 + 5\eu^{-6.651} = 1.0065$; विषम योग $3(3\eu^{-2.217} + 7\eu^{-13.30}) = 0.9806$;
$x_{\mathrm{para}} = 1.0065/1.9871 = 0.51$।
\end{solution}

\begin{solution}{exo:b3:statistical-thermo-applied:4}
$I = 1$ के लिए $3 \times 2 = 6$ सममित और $3 \times 1 = 3$ प्रतिसममित चक्रण अवस्थाएँ हैं;
ड्यूटेरॉन बोसॉन हैं, अतः पूर्ण तरंग फलन सममित है: सममित चक्रण
अवस्थाएँ सम $J$ के साथ (ऑर्थो, भार 6), प्रतिसममित विषम $J$ के साथ (पैरा, भार
3)। ऑर्थो-ड्यूटेरियम, जिसमें $J = 0$ है, 0 K पर स्थायी है; कक्ष ताप पर
ऑर्थो:पैरा = 2:1।
\end{solution}

\begin{solution}{exo:b3:statistical-thermo-applied:5}
$\Delta_rE_0 = 2 \times 1883.7 - 2170.3 - 1542.3 = \qty{54.8}{cm^{-1}}$; $\eu^{-54.8/207.2} = 0.768$, और
$4 \times 0.768 = 3.07$। पूर्ण मान, 3.26, 6\,\% अधिक है: स्थानांतरणीय, घूर्णी और
कंपनिक गुणक केवल चिरसम्मत सीमा में यथार्थतः निरस्त होते हैं, और
इन हल्के अणुओं का घूर्णन \qty{298}{K} पर अब भी आंशिक रूप से क्वांटीकृत है।
\end{solution}

\begin{solution}{exo:b3:statistical-thermo-applied:6}
\ce{N2O}: दो अभिविन्यास, $R\ln2 = \qty{5.76}{J.K^{-1}.mol^{-1}}$। \ce{CH3D}: चार, $R\ln4 =
\qty{11.53}{J.K^{-1}.mol^{-1}}$ (उपरि सीमाएँ, जो व्यवस्थाएँ यादृच्छिक होने पर प्राप्त होती हैं)।
\end{solution}

\begin{solution}{exo:b3:statistical-thermo-applied:7}
$x = 797.6/300 = 2.659$: आइंस्टाइन फलन $x^2\eu^{-x}/(1 - \eu^{-x})^2 = 0.572$; $C_V/R = 3.072$,
$C_p = 4.072R = \qty{33.86}{J.K^{-1}.mol^{-1}}$, सारणी से 0.4\,\% नीचे (अनावर्तिता)।
\end{solution}

\begin{solution}{exo:b3:statistical-thermo-applied:8}
$x = 1.029$: $x^2\eu^{-x}/(1 - \eu^{-x})^2 = 0.916$, अर्थात् \qty{7.62}{J.K^{-1}.mol^{-1}}, $R$ का 92\,\%:
कक्ष ताप पर आयोडीन का कंपन लगभग चिरसम्मत है।
\end{solution}

\begin{solution}{exo:b3:statistical-thermo-applied:9}
B, \ce{^{18}O} में अधिक समृद्ध है। $\alpha_{\mathrm{A/B}} = (1 - 0.0100)/(1 + 0.0020) = 0.98802$।
\end{solution}

\begin{solution}{exo:b3:statistical-thermo-applied:10}
$\Delta\mathrm{ZPE} = \frac12(1 - 1/\sqrt2)(2000 - 3000) = \qty{-146}{cm^{-1}}$; $K = \eu^{146.4/207.2} = 2.0$।
ड्यूटेरियम वरीयता से X, अर्थात् अधिक दृढ़ आबंध, पर जाता है, जहाँ उसकी \omterm{def:b3:quantum-model-systems:oscillator}{शून्य-बिंदु
ऊर्जा} की बचत सबसे अधिक है।
\end{solution}

\begin{solution}{exo:b3:statistical-thermo-applied:11}
$\Delta_rH^\circ = R\ln10\,(-1.766 + 3.392)/(1/900 - 1/1100) = 19.145 \times 1.626/\num{2.020e-4} =
\qty{154}{kJ/mol}$। दोनों परमाणु $2 \times \frac52RT$ एन्थैल्पी वहन करते हैं; अणु $\frac52RT + RT$
(घूर्णन) और अपनी कंपनिक ऊर्जा $R\theta_{\mathrm V}/(\eu^{\theta_{\mathrm V}/T} - 1)$, जो $RT$ से थोड़ी कम है:
अंतर, \qty{1000}{K} पर लगभग \qty{5}{kJ/mol}, $D_0$ में जुड़ता है।
\end{solution}

\begin{solution}{exo:b3:statistical-thermo-applied:12}
स्तर 0 ($g = 1$) और $6hcB$ ($g = 5$): $q = 1 + 5\eu^{-x}$, $\langle\varepsilon\rangle/kT = 5x\eu^{-x}/q$, $\langle\varepsilon^2\rangle/(kT)^2 =
5x^2\eu^{-x}/q$, और $C/R = \langle\varepsilon^2\rangle/(kT)^2 - (\langle\varepsilon\rangle/kT)^2 = 5x^2\eu^{-x}/(1 + 5\eu^{-x})^2$।
\qty{100}{K} पर $x = 5.121$: $C/R = 0.783/1.061 = 0.738$।
\end{solution}

\begin{solution}{pb:b3:statistical-thermo-applied:1}
\textbf{1.} $2 \times 2 = 4$: तीन सममित (त्रिक), एक प्रतिसममित (एकल)।
\textbf{2.} प्रोटॉन फ़र्मिऑन हैं, अतः विनिमय पर पूर्ण तरंग फलन चिह्न बदलता है;
घूर्णी फलन $(-1)^J$ देता है: सम $J$
प्रतिसममित एकल (पैरा) के साथ जाता है, विषम $J$ सममित त्रिक (ऑर्थो) के साथ।
\textbf{3.} सम स्तर 1, विषम स्तर 3।
\textbf{4.} \qty{300}{K} ($T = 3.5\,\theta_{\mathrm R}$) पर सम और विषम घूर्णी योग लगभग बराबर हैं;
1 और 3 से भारित होकर वे 1:3 देते हैं (पैरा अंश 0.2507)।
\textbf{5.} $F(1) = 2 \times 59.322 - 4 \times 0.0471 = \qty{118.46}{cm^{-1}} = \qty{1.417}{kJ/mol}$।
\textbf{6.} $1/(1 + 9\eu^{-170.4/20.37}) = 1/(1 + 9 \times \num{2.3e-4}) = 0.998$।
\textbf{7.} केवल $J = 0$ और 1 के साथ, आधे पैरा का अर्थ है $9\eu^{-2\theta_{\mathrm R}/T} = 1$, $T = 2\theta_{\mathrm R}/\ln9 =
0.91\,\theta_{\mathrm R} = \qty{78}{K}$: प्रतिस्पर्धा $J = 1$ के भार 9 और उसके
बोल्ट्समान गुणक के बीच है।
\textbf{8.} 0.75। उत्प्रेरक के बिना रूपांतरण को दिन लगते हैं, द्रवण को
घंटे।
\textbf{9.} $(0.75 - 0.002) \times 1.417 = \qty{1.060}{kJ/mol}$।
\textbf{10.} $1.060/\num{2.016e-3} = \qty{526}{kJ/kg}$।
\textbf{11.} $0.905/\num{2.016e-3} = \qty{449}{kJ/kg}$।
\textbf{12.} रूपांतरण की ऊष्मा वाष्पन एन्थैल्पी की 1.17 गुनी है।
\textbf{13.} यह द्रव के भीतर उत्पन्न होती है, दीवारों से होकर नहीं आती।
\textbf{14.} $x_{\mathrm o}(t) = x_{\mathrm{eq}} + (0.75 - x_{\mathrm{eq}})\eu^{-kt}$, जहाँ $x_{\mathrm{eq}} = 0.002$।
\textbf{15.} $0.002 + 0.748\eu^{-0.274} = 0.571$।
\textbf{16.} $(0.75 - 0.571) \times 1.417 = \qty{0.254}{kJ/mol}$।
\textbf{17.} $0.254/0.905 = 0.28$: एक दिन में टंकी का 28\,\%।
\textbf{18.} \qty{168}{h} के बाद $x_{\mathrm o} = 0.112$, ऊष्मा \qty{0.904}{kJ/mol}, जो वाष्पन
एन्थैल्पी के बराबर है: इस मॉडल में टंकी खाली हो जाती है। आकलन हानि को अधिक बताता है
क्योंकि वाष्पित अणु अपना अरूपांतरित \omterm{def:b3:statistical-thermo-applied:spin-isomers}{ऑर्थो हाइड्रोजन} साथ ले
जाते हैं; वास्तविक हानि कम है परंतु फिर भी बड़ी।
\textbf{19.} $t_{1/2} = \ln2/0.0114 = \qty{61}{h}$।
\textbf{20.} द्रवित्र में, क्रमिक तापों पर उत्प्रेरक संस्तरों पर, ताकि
प्रशीतक द्रव के भंडारण से पहले रूपांतरण ऊष्मा हटा दे।
\textbf{21.} हाँ: \qty{300}{K} पर साम्य पैरा अंश 0.2507 है।
\textbf{22.} सामान्य हाइड्रोजन में दोनों रूप परस्पर नहीं बदल सकते, अतः हर एक का
घूर्णन अपनी स्तर-सीढ़ी पर अलग जमता है; साम्य हाइड्रोजन में
दी गई ऊष्मा का कुछ भाग पैरा को ऑर्थो में बदलता है।
\textbf{23.} लगभग \qty{50}{K} के निकट लगभग 3.57: घूर्णन उत्तेजित करने के अतिरिक्त ऊष्मा
अणुओं को $J = 0$ से $J = 1$ में, \qty{118}{cm^{-1}} ऊपर, बदलती है, जो धनात्मक
एन्थैल्पी वाली एक अभिक्रिया है।
\textbf{24.} 1.50: दोनों रूपों में घूर्णन जमा हुआ है (पैरा $J = 0$ में, ऑर्थो $J = 1$ में)।
\textbf{25.} \textbf{सामान्य हाइड्रोजन के लिए रूपांतरण ऊष्मा / वाष्पन एन्थैल्पी $\approx 1.2$}:
अकेला उसका रूपांतरण पूरी टंकी को वाष्पित कर सकता है, इसीलिए
भंडारण से पहले हाइड्रोजन को पैरा में बदला जाता है।
\end{solution}
